\begin{proof}[Proof of \cref{obj:lem:finite-prefix-transfer}]
\begin{enumerate}
\item For a length vector \(\mathbf j=(j_i)_{i\in I}\in\mathbb N_0^I\), define
\begingroup\small
\[
\omega(\mathbf j)
=
\prod_{i\in I}
e^{-h_i/8}\frac{(h_i/8)^{j_i}}{j_i!},
\qquad
\mathsf{Good}
=
\{(s,\mathbf j):j_i\le h_i\text{ for every }i\in I\},
\qquad
\mathsf{Overflow}=\mathsf{Good}^{\,c}.
\]
\endgroup
Here \(s\) ranges over the finite ordered pool families in the statement. Define, for every such \(s\) and every \(\mathbf j\in\mathbb N_0^I\),
\[
T_{\mathrm{cap}}(s,\mathbf j)
=
\begin{cases}
T\bigl((s_{i,1:j_i})_{i\in I}\bigr),
   &(s,\mathbf j)\in\mathsf{Good},\\
0,&(s,\mathbf j)\in\mathsf{Overflow}.
\end{cases}
\]
On \(\mathsf{Good}\), the assumed range of \(T\) gives
\(\lvert T_{\mathrm{cap}}(s,\mathbf j)\rvert\le1\); on
\(\mathsf{Overflow}\), this follows from its zero value. Thus
\[
\lvert T_{\mathrm{cap}}(s,\mathbf j)\rvert\le1
\quad\text{for every }(s,\mathbf j).
\]
Writing \(\mathbb E_J\) for expectation over the independent prefix lengths, we have
\[
\begin{aligned}
\mathbb E_J T_{\mathrm{cap}}(s,J)
&=
\sum_{\mathbf j\in\mathbb N_0^I}
\omega(\mathbf j)T_{\mathrm{cap}}(s,\mathbf j)\\
&=
\sum_{\substack{\mathbf j\in\mathbb N_0^I\\j_i\le h_i\ \forall i}}
\omega(\mathbf j)
T\bigl((s_{i,1:j_i})_{i\in I}\bigr)
=
T_{\mathrm{fin}}(s).
\end{aligned}
\]
Every prefix projection in the last sum is measurable, and the sum is finite. Consequently \(T_{\mathrm{fin}}\) is a total deterministic measurable rule, with
\[
\lvert T_{\mathrm{fin}}(s)\rvert
\le
\mathbb E_J\lvert T_{\mathrm{cap}}(s,J)\rvert
\le1.
\]
For Borel observation spaces, these compositions and their finite sum are Borel measurable. % lean: Causalean.Stat.FiniteRaoBlackwell.IndependentPoissonPrefix.FiniteFamily.finite_prefix_transfer

\item Write \(\mathbb E_s\) for expectation under the independent finite iid pools and
\[
\mathbb E_{s,J}[\cdot]=\mathbb E_s\mathbb E_J[\cdot].
\]
The bounds on \(T_{\mathrm{cap}}\) and \(\vartheta\) imply
\[
\lvert T_{\mathrm{cap}}(s,J)-\vartheta\rvert\le2,
\qquad
0\le(T_{\mathrm{cap}}(s,J)-\vartheta)^2\le4,
\]
so all the following losses are integrable. For each fixed \(s\), Jensen's inequality for the square function gives
\[
\begin{aligned}
(T_{\mathrm{fin}}(s)-\vartheta)^2
&=
\left(\mathbb E_J[T_{\mathrm{cap}}(s,J)-\vartheta]\right)^2\\
&\le
\mathbb E_J[(T_{\mathrm{cap}}(s,J)-\vartheta)^2].
\end{aligned}
\]
Integrating over the finite pools yields
\[
\mathbb E_s[(T_{\mathrm{fin}}(s)-\vartheta)^2]
\le
\mathbb E_{s,J}[(T_{\mathrm{cap}}(s,J)-\vartheta)^2].
\] % lean: Causalean.Stat.FiniteRaoBlackwell.IndependentPoissonPrefix.FiniteFamily.sqRisk_prefixRaoBlackwellStatistic_le_capped

\item Let \(\boldsymbol s^{\mathrm{Poi}}\) denote the untruncated random prefix family from the hypothesis. Its conditional law is
\[
\mathcal L(\boldsymbol s^{\mathrm{Poi}}\mid J=\mathbf j)
=
\bigotimes_{i\in I}P_i^{\otimes j_i},
\qquad
\mathbf j\in\mathbb N_0^I,
\]
where
\[
P_i^{\otimes0}=\delta_{()},
\]
the point mass on the empty ordered prefix. For every admissible \(\mathbf j\), independence and iid sampling give the same conditional prefix law for the finite pools:
\[
\mathcal L\bigl((s_{i,1:j_i})_{i\in I}\mid J=\mathbf j\bigr)
=
\bigotimes_{i\in I}P_i^{\otimes j_i},
\qquad j_i\le h_i\ \forall i.
\]
Indeed, projection onto the first \(j_i\) coordinates sends
\(P_i^{\otimes h_i}\) to \(P_i^{\otimes j_i}\), and the coordinatewise laws multiply across the independent pools. Summing these equalities over the admissible lengths therefore gives
\[
\begin{aligned}
&\mathbb E_{s,J}\!\left[
(T_{\mathrm{cap}}(s,J)-\vartheta)^2
\mathbf1_{\mathsf{Good}}(s,J)\right]\\
&\quad=
\sum_{\substack{\mathbf j\in\mathbb N_0^I\\j_i\le h_i\ \forall i}}
\omega(\mathbf j)
\int
(T(s')-\vartheta)^2\,
d\!\left(\bigotimes_{i\in I}P_i^{\otimes j_i}\right)(s')\\
&\quad=
\mathbb E\!\left[
(T(\boldsymbol s^{\mathrm{Poi}})-\vartheta)^2
\mathbf1_{\{J_i\le h_i\ \forall i\}}\right].
\end{aligned}
\]
In the integral, \(s'\) ranges over prefix families with respective lengths \(j_i\). % lean: Causalean.Stat.FiniteRaoBlackwell.IndependentPoissonPrefix.FiniteFamily.sqRisk_cappedPrefixStatistic_restrict_nonoverflow_eq

Splitting the capped loss between the two events, and using its zero output on overflow, now yields
\[
\begin{aligned}
\mathbb E_{s,J}[(T_{\mathrm{cap}}(s,J)-\vartheta)^2]
&=
\mathbb E\!\left[
(T(\boldsymbol s^{\mathrm{Poi}})-\vartheta)^2
\mathbf1_{\{J_i\le h_i\ \forall i\}}\right]\\
&\qquad+
\vartheta^2\Pr(\exists i\in I:J_i>h_i)\\
&\le
\mathbb E[(T(\boldsymbol s^{\mathrm{Poi}})-\vartheta)^2]
+\Pr(\exists i\in I:J_i>h_i)\\
&\le
\mathbb E[(T(\boldsymbol s^{\mathrm{Poi}})-\vartheta)^2]
+\sum_{i\in I}\Pr(J_i>h_i).
\end{aligned}
\]
The first inequality discards a nonnegative complementary loss and uses
\(\vartheta^2\le1\). The last follows by writing the overflow event as
\(\bigcup_{i\in I}\{J_i>h_i\}\) and applying the finite union bound. % lean: Causalean.Stat.FiniteRaoBlackwell.IndependentPoissonPrefix.FiniteFamily.zero_overflow_prefix_risk

\item Fix \(i\in I\). The Poisson exponential moment follows directly from its probability masses and the exponential series:
\[
\begin{aligned}
\mathbb E[8^{J_i}]
&=
e^{-h_i/8}
\sum_{k=0}^{\infty}
\frac{(h_i/8)^k8^k}{k!}\\
&=
e^{-h_i/8}
\sum_{k=0}^{\infty}\frac{h_i^k}{k!}
=
e^{7h_i/8}.
\end{aligned}
\]
Since \(\{J_i>h_i\}\subseteq\{J_i\ge h_i\}\), exponential Markov gives
\[
\Pr(J_i>h_i)e^{-7h_i/8}
\le
e^{-h_i\log8}.
\]
The elementary numerical lower bound \(\log2>5/8\) implies
\[
\log8=3\log2\ge\frac{15}{8}.
\]
As \(h_i\ge0\), we obtain
\[
\Pr(J_i>h_i)e^{-7h_i/8}
\le
e^{-h_i\log8}
\le
e^{-15h_i/8}
=
e^{-h_i}e^{-7h_i/8}.
\]
Cancelling the strictly positive factor \(e^{-7h_i/8}\) proves
\[
\Pr(J_i>h_i)\le e^{-h_i}.
\] % lean: Causalean.Stat.FiniteRaoBlackwell.IndependentPoissonPrefix.FiniteFamily.eighth_pool_poisson_overflow

\item Combining the preceding bounds with the assumed Poisson-prefix risk gives
\[
\begin{aligned}
\mathbb E_s[(T_{\mathrm{fin}}(s)-\vartheta)^2]
&\le
\mathbb E[(T(\boldsymbol s^{\mathrm{Poi}})-\vartheta)^2]
+\sum_{i\in I}\Pr(J_i>h_i)\\
&\le
A+\sum_{i\in I}e^{-h_i}.
\end{aligned}
\]
These calculations include \(I=\varnothing\): empty products equal \(1\), empty sums equal \(0\), the pool and prefix families are the unique empty families, and the overflow event is empty. % lean: CausalSmith.Stat.AnnotationRarearmFrontier.finite_prefix_transfer
\end{enumerate}
\end{proof}
